% =====================================================================
%  1.3 Rotuttrykk og brøkeksponenter — ekstra utfordring
% =====================================================================

\begin{infoboks}{Velg metode selv}
Arket bygger på 1.1 Potenser, 1.2 Kvadratsetningene og 1.3 Rotuttrykk. Vis regning og begrunnelse.
Bruk eksakte, forenklede svar. Kontroller svaret ved å sette inn tall.
\end{infoboks}

\nivaahode{2}

\oppg{U1}{2}{\uh}
Konjugatsetningen fra 1.2 gir $(\sqrt{3}-1)(\sqrt{3}+1) = 3 - 1 = 2$. Utvid brøken med konjugatet og gjør nevneren rasjonal.
\begin{tasks}(3)
  \task $\dfrac{1}{\sqrt{3}-1}$
  \task $\dfrac{1}{\sqrt{5}+2}$
  \task $\dfrac{6}{\sqrt{7}-1}$
\end{tasks}

\oppg{U2}{2}{\uh}
Gang ut og forenkle.
\begin{tasks}(3)
  \task $(2+\sqrt{3})^2$
  \task $(\sqrt{5}+\sqrt{2})(\sqrt{5}-\sqrt{2})$
  \task $(\sqrt{3}+\sqrt{2})^2 - (\sqrt{3}-\sqrt{2})^2$
\end{tasks}

\oppg{U3}{2}{\uh}
Løs likningene. Skriv begge sider med samme grunntall først.
\begin{tasks}(3)
  \task $2^x = \sqrt{8}$
  \task $9^x = \sqrt[3]{3}$
  \task $8^x = \dfrac{1}{\sqrt{2}}$
\end{tasks}

\nivaahode{3}

\oppg{U4}{3}{\prog}
Herons metode finner $\sqrt{2}$ ved å gjenta formelen under, med start i $x = 1$.
\[ x \to \dfrac{x + \dfrac{2}{x}}{2} \]
\begin{kode}
x = 1
runder = 0
while abs(x*x - 2) > 1e-10:
    x = (x + 2/x)/2
    runder = runder + 1
print(runder)
\end{kode}
\begin{tasks}(1)
  \task Regn for hånd verdien av $x$ etter runde 1 og runde 2. Svar som brøker.
  \task Kjør programmet. Hvor mange runder trengs?
  \task Hvis $x$ er for stor, er $\frac{2}{x}$ for liten. Forklar hvorfor gjennomsnittet av dem kommer nærmere $\sqrt{2}$.
  \task Endre programmet slik at det finner $\sqrt{10}$. Hvor mange runder trengs nå?
\end{tasks}

\oppg{U5}{3}{\uh}
Ordne $\sqrt{2}$, $\sqrt[3]{3}$ og $\sqrt[6]{6}$ etter størrelse, minst først. Begrunn uten kalkulator.
\hint{Hvilken potens fjerner alle rottegnene?}

\oppg{U6}{3}{\uh}
Vi skal vise at $\sqrt{2}$ er irrasjonalt. Anta det motsatte: $\sqrt{2} = \dfrac{p}{q}$, der $p$ og $q$ er heltall, $q \neq 0$, og de har ingen felles faktor større enn 1.
\begin{tasks}(1)
  \task Kvadrer begge sider og fullfør: $p^2 = \skrivlinje[2cm]\cdot q^2$.
  \task Da er $p^2$ et partall. Forklar hvorfor $p$ da må være et partall.
  \task Sett $p = 2k$. Vis at $q^2 = \skrivlinje[2cm]$.
  \task Forklar hvorfor $q$ også er et partall, og hvorfor det gir en selvmotsigelse.
  \task Prøv de samme stegene med $\sqrt{4} = \dfrac{p}{q}$. Hvilket steg bryter sammen, og hvorfor?
\end{tasks}

\oppg{U7}{3}{\uh}
Lag selv et uttrykk på formen $\dfrac{c}{\sqrt{a}+\sqrt{b}}$ som blir et uttrykk uten brøk når nevneren gjøres rasjonal ved å utvide med konjugatet. Vis utregningen. Hvilken betingelse må $a$, $b$ og $c$ oppfylle?

\begin{oppsummering}
En rot er en potens med brøkeksponent. Rottegnet gjelder produkt og brøk, men ikke sum.
Gjør nevneren rasjonal ved å utvide, eventuelt med konjugatet.
Kontroller ved å sette inn tall eller gange tilbake.
\end{oppsummering}
